\chapter{Il modulo \code{sync-queue}}
\label{cap:10-sync-queue}

Il modulo \code{sync-queue} inserisce sul percorso sincrono una coda esplicita
con controllo di ammissione. \`E l'applicazione diretta del principio di
SEDA~\cite{welsh2001seda}: il punto in cui le richieste si accumulano dev'essere
noto e strumentato, anzich\'e disperso nei buffer del sistema operativo o nella
lunghezza delle connessioni pendenti.

Il modulo risponde a una domanda che il \code{async-queue} non pone: \emph{questa
richiesta ha senso accettarla?} Una coda senza ammissione accetta finch\'e ha
posto e rifiuta quando \`e piena, il che significa che l'ultima richiesta
accettata prima del riempimento attende quanto l'intera coda impiega a
drenare. Con l'ammissione, il rifiuto avviene prima, sulla base di una stima.

Sono 905 righe in dieci classi.

\section{Le due domande dell'ammissione}

\code{SyncQueueAdmissionController} \`e la classe pi\`u corta e pi\`u
significativa del modulo:

\begin{lstlisting}[language=Java]
public SyncQueueAdmissionResult evaluate(String functionName, int depth, Instant now) {
    if (depth >= maxDepth) {
        return SyncQueueAdmissionResult.rejected(SyncQueueRejectReason.DEPTH,
                                                 Double.POSITIVE_INFINITY);
    }
    double estWaitSeconds = estimator.estimateWaitSeconds(functionName, depth, now);
    long maxWaitSeconds = configSource.syncQueueMaxEstimatedWait().toSeconds();
    if (configSource.syncQueueAdmissionEnabled()
            && (maxWaitSeconds == 0 || estWaitSeconds > maxWaitSeconds)) {
        return SyncQueueAdmissionResult.rejected(SyncQueueRejectReason.EST_WAIT, estWaitSeconds);
    }
    return SyncQueueAdmissionResult.accepted(estWaitSeconds);
}
\end{lstlisting}

Le due condizioni sono qualitativamente diverse.

\paragraph{\code{DEPTH}} \`e un vincolo di risorsa. La coda ha una capacit\`a
finita e non pu\`o accettare oltre. Il rifiuto \`e incondizionato e non dipende
dalla configurazione dell'ammissione: anche disattivando il controllo di
ammissione la coda continua a rifiutare quando \`e piena.

\paragraph{\code{EST\_WAIT}} \`e un vincolo di qualit\`a del servizio. La coda
avrebbe posto, ma la richiesta attenderebbe troppo a lungo. Il rifiuto \`e una
scelta di politica, disattivabile.

La distinzione \`e resa visibile al chiamante attraverso l'header
\code{X-Queue-Reject-Reason}, ed \`e la stessa distinzione su cui il modulo
\code{offload} basa la propria attivazione reattiva
(\cref{cap:13-offload}): entrambe le condizioni sono \emph{pressione}, e
giustificano una deviazione verso un'istanza remota, mentre un
\code{TIMEOUT} in coda non lo \`e.

\section{La stima dell'attesa}

\subsection{Il modello}

\code{WaitEstimator} applica la forma pi\`u diretta della legge di
Little~\cite{little1961proof}. Se il sistema smaltisce $\mu$ richieste al secondo
e ne ha $Q$ in coda, una richiesta che si accodasse adesso attenderebbe
approssimativamente

\[
  W_{\text{stim}} = \frac{Q}{\mu}.
\]

Il throughput $\mu$ \`e stimato contando i dispatch effettivi in una finestra
scorrevole:

\begin{lstlisting}[language=Java]
public double estimateWaitSeconds(String functionName, int queueDepth, Instant now) {
    if (queueDepth <= 0) { return 0.0; }
    ThroughputSnapshot perFunction = snapshot(perFunctionEvents.get(functionName), now);
    if (perFunction.samples() >= perFunctionMinSamples && perFunction.throughput() > 0) {
        return queueDepth / perFunction.throughput();
    }
    ThroughputSnapshot global = snapshot(globalEvents, now);
    if (global.throughput() <= 0) { return Double.POSITIVE_INFINITY; }
    return queueDepth / global.throughput();
}
\end{lstlisting}

\subsection{Le tre decisioni della stima}

\paragraph{Prima: per funzione, se ci sono abbastanza campioni.} Funzioni diverse
hanno tempi di servizio diversi, quindi la stima per funzione \`e la sola
significativa. Ma una stima costruita su pochi campioni \`e rumorosa, e un
rumore su questa grandezza si traduce direttamente in rifiuti ingiustificati. Il
parametro \code{perFunctionMinSamples} (default 50) fissa la soglia sotto la
quale la stima per funzione non \`e considerata affidabile.

\paragraph{Seconda: altrimenti, globale.} Sotto la soglia si ripiega sul
throughput aggregato della piattaforma. \`E una stima peggiore --- mescola
funzioni con costi diversi --- ma \`e disponibile prima. Il ripiego \`e una
scelta pragmatica: preferire una stima grossolana a nessuna stima.

\paragraph{Terza: senza campioni, attesa infinita.} \`E il caso pi\`u
interessante e ha una conseguenza operativa che il progetto ha appreso per via
empirica. Con \code{globalThroughput = 0} la stima \`e
\code{Double.POSITIVE\_INFINITY}, che eccede qualunque soglia, e quindi
\emph{ogni} richiesta viene rifiutata con \code{429}.

Questo \`e corretto in produzione --- un sistema che non ha ancora dispacciato
nulla non pu\`o promettere alcun tempo di attesa --- e distruttivo nei test: una
suite che avvia il control plane e invia la prima richiesta la vede respinta con
\code{429} senza alcuna congestione reale. La soluzione adottata \`e che i test
di API disattivano esplicitamente la coda sincrona
(\code{sync-queue.enabled=false}), il che \`e registrato come vincolo noto
dell'impianto di test (\cref{cap:26-qualita-verifica}).

\`E anche un esempio del genere di comportamento che una piattaforma di ricerca
deve documentare: la condizione ``nessuna storia'' non \`e un caso limite raro,
\`e lo stato in cui il sistema si trova ogni volta che viene avviato.

\subsection{La finestra}

Il default \`e una finestra di 30 secondi. Il calcolo del throughput divide il
numero di campioni per la durata della finestra, non per l'intervallo effettivo
fra il primo e l'ultimo campione:

\begin{lstlisting}[language=Java]
double seconds = Math.max(1.0, window.toSeconds());
return new ThroughputSnapshot(samples, samples / seconds);
\end{lstlisting}

Questo \emph{sottostima} il throughput durante i primi 30 secondi di attivit\`a,
perch\'e divide per una finestra che il sistema non ha ancora riempito. Una
sottostima del throughput produce una sovrastima dell'attesa, e quindi un
comportamento pi\`u conservativo --- pi\`u rifiuti --- nella fase iniziale. La
direzione dell'errore \`e quindi sicura, ed \`e presumibilmente il motivo per cui
la formulazione \`e stata mantenuta; il progetto non lo documenta esplicitamente.

\section{Lo scheduler sincrono e il forward progress}

Il consumatore della coda \`e \code{SyncScheduler}, e la sua funzione
\code{tickOnce()} risolve un problema che una coda FIFO ingenua non risolve.

\subsection{Il problema della testa bloccata}

Con una singola coda condivisa da tutte le funzioni, la richiesta in testa
appartiene a una funzione che potrebbe non avere slot di concorrenza liberi. Un
consumatore FIFO stretto si bloccherebbe su di essa, e nessuna richiesta di
nessun'altra funzione verrebbe dispacciata --- \emph{head-of-line blocking}, con
la variante particolarmente sgradevole che il blocco dura quanto il tempo di
servizio della funzione bloccante.

\subsection{La soluzione: scansione selettiva}

\begin{lstlisting}[language=Java, caption={Il cuore di \code{tickOnce()}.}]
SyncQueueItem item = queue.findReadyMatching(now,
        task -> enqueuer.hasAvailableSlot(task.functionName()));
if (item == null) {
    if (queue.peekReady(now) == null) {
        blockedBackoffMs = tickMs;
        queue.awaitWork(tickMs);          // coda vuota: attendi lavoro
    } else {
        queue.rotateReadyScanWindow(now); // tutti bloccati: ruota e ritenta
        pause.accept(currentBlockedBackoff());
    }
    return;
}
String functionName = item.task().functionName();
if (!enqueuer.tryAcquireSlot(functionName)) {
    queue.rotateReadyItem(item, now);     // slot perso in corsa: rimetti in fondo
    pause.accept(currentBlockedBackoff());
    return;
}
...
\end{lstlisting}

Lo scheduler non prende la testa: \emph{cerca} il primo elemento la cui funzione
ha uno slot disponibile. La ricerca \`e limitata a una finestra di 64 elementi
(\code{POLL\_READY\_MATCHING\_SCAN\_LIMIT}), il che mantiene il costo per tick
limitato indipendentemente dalla profondit\`a della coda.

Se nessun elemento nella finestra \`e servibile, la finestra viene \emph{ruotata}
--- i suoi elementi vanno in fondo --- cos\`i che il tick successivo esamini
elementi diversi. La rotazione \`e ci\`o che rende la scansione limitata
comunque equa: nessun elemento resta indefinitamente fuori dalla finestra
esaminata.

Il codice tratta anche il caso di corsa: fra il momento in cui un elemento \`e
selezionato perch\'e la sua funzione aveva uno slot e il momento in cui lo slot
viene acquisito, un altro consumatore pu\`o averlo preso. In tal caso l'elemento
viene ruotato anzich\'e scartato.

\subsection{Il backoff}

\begin{lstlisting}[language=Java]
private volatile long tickMs = 2;
private volatile long blockedBackoffMs = tickMs;

private long currentBlockedBackoff() {
    long current = blockedBackoffMs;
    blockedBackoffMs = Math.min(current * 2, 50);
    return current;
}
\end{lstlisting}

Il tick base \`e di 2\,ms. Quando lo scheduler trova lavoro ma non riesce a
dispacciarlo --- ossia quando ogni funzione in coda \`e satura --- attende con
backoff esponenziale fino a un tetto di 50\,ms. Il backoff viene azzerato appena
un dispatch riesce.

I due valori sono un compromesso esplicito fra reattivit\`a e consumo di CPU. A
2\,ms il ciclo aggiunge al pi\`u due millisecondi alla latenza di una richiesta
che attende uno slot appena liberato, il che \`e dello stesso ordine del tempo di
servizio delle funzioni di riferimento del progetto. Senza backoff, un sistema
saturo farebbe girare il thread dello scheduler a vuoto 500 volte al secondo.

\begin{damisurare}
Il contributo del tick di 2\,ms alla latenza end-to-end non \`e stato isolato
sperimentalmente. Il \cref{cap:25-valutazione-concorrenza} riporta tempi di
attesa dell'ordine di 40\,ms contro tempi di servizio di 5\,ms, in cui il
contributo del tick sarebbe minoritario ma non trascurabile.
\todomis{contributo del periodo di tick alla latenza a bassa concorrenza}
\end{damisurare}

\section{Timeout in coda}

Un elemento pu\`o scadere mentre attende. La scadenza \`e governata da
\code{sync-queue.max-queue-wait} (default 2\,s) e viene verificata a ogni
ispezione della coda: \code{peekReady}, \code{findReadyMatching} e le rotazioni
rimuovono gli elementi scaduti e li completano con esito \code{QUEUE\_TIMEOUT}.

Il coordinatore reattivo traduce quell'esito in un rifiuto con ragione
\code{TIMEOUT}:

\begin{lstlisting}[language=Java]
if (result.error() != null && "QUEUE_TIMEOUT".equals(result.error().code())) {
    throw new SyncQueueRejectedException(SyncQueueRejectReason.TIMEOUT,
                                         syncQueueGateway.retryAfterSeconds());
}
\end{lstlisting}

Il chiamante riceve quindi \code{429} con \code{X-Queue-Reject-Reason: timeout},
distinguibile dai due rifiuti di ammissione. La distinzione conta: un rifiuto per
\code{DEPTH} o \code{EST\_WAIT} avviene \emph{prima} che la richiesta occupi la
coda, un \code{TIMEOUT} avviene dopo che l'ha occupata per l'intera durata
consentita. Solo i primi due sono trattati come pressione dal modulo di offload.

\section{Rimozione di una funzione}

Quando una funzione viene cancellata mentre ha richieste in coda, il modulo
esegue tre azioni: registra il nome in un insieme di funzioni rimosse, drena
dalla coda tutti gli elementi che le appartengono completandoli con errore
\code{FUNCTION\_REMOVED}, e cancella lo stato dell'estimatore e delle metriche.

L'insieme delle funzioni rimosse \`e consultato anche in ingresso, due volte: una
prima del lock sulla coda e una dentro. La doppia verifica gestisce la corsa fra
una registrazione in corso e una cancellazione, e ha il costo di una lettura su
un insieme concorrente, trascurabile rispetto alla sezione critica che protegge.

\section{Metriche esportate}

\begin{center}
\small
\begin{adjustbox}{max width=\textwidth}
\begin{tabular}{@{}lll@{}}
\toprule
\textbf{Metrica} & \textbf{Tipo} & \textbf{Significato} \\
\midrule
\code{sync\_queue\_depth} & gauge & Profondit\`a, globale (\code{function=""}) e per funzione. \\
\code{sync\_queue\_wait\_seconds} & gauge & Attesa stimata, globale e per funzione. \\
\code{sync\_queue\_admitted\_total} & counter & Ammissioni. \\
\code{sync\_queue\_rejected\_total} & counter & Rifiuti di ammissione. \\
\code{sync\_queue\_timedout\_total} & counter & Scadenze in coda. \\
\bottomrule
\end{tabular}
\end{adjustbox}
\end{center}

La documentazione operativa del progetto fornisce una regola di lettura che vale
la pena riportare perch\'e sintetizza il rapporto fra le due grandezze:
\emph{``un valore alto di \code{sync\_queue\_rejected\_total\{function\}} con
profondit\`a bassa indica un rifiuto per attesa stimata, non un esaurimento della
capacit\`a''}. \`E il modo in cui un operatore distingue una piattaforma che ha
finito la coda da una che sta proteggendo il proprio SLO.

\section{Configurazione}

\begin{center}
\small
\begin{adjustbox}{max width=\textwidth}
\begin{tabular}{@{}llp{6.4cm}@{}}
\toprule
\textbf{Chiave} & \textbf{Default} & \textbf{Effetto} \\
\midrule
\code{enabled} & \code{true} & Attiva il modulo. \\
\code{admission-enabled} & \code{true} & Attiva il rifiuto per \code{EST\_WAIT}
(quello per \code{DEPTH} resta). \\
\code{max-depth} & 200 & Capacit\`a della coda. \\
\code{max-estimated-wait} & 2\,s & Soglia di rifiuto per attesa stimata. \\
\code{max-queue-wait} & 2\,s & Scadenza in coda. \\
\code{retry-after-seconds} & 2 & Valore dell'header \code{Retry-After}. \\
\code{throughput-window} & 30\,s & Finestra della stima di throughput. \\
\code{per-function-min-samples} & 50 & Campioni per fidarsi della stima per
funzione. \\
\bottomrule
\end{tabular}
\end{adjustbox}
\end{center}

Tutti questi parametri, tranne \code{max-depth} e i due che governano la stima,
sono modificabili a caldo attraverso il modulo \code{runtime-config}
(\cref{cap:14-runtime-config}): la lettura avviene infatti attraverso
\code{SyncQueueConfigSource} anzich\'e direttamente dalle propriet\`a. \`E il
motivo per cui \code{evaluate()} interroga il config source a ogni valutazione
anzich\'e leggere un campo finale.
